\chapter{Metodologia}
\label{cap:metodologia}

De acordo com os objetivos dispostos no Capítulo 1, a metodologia deste trabalho estrutura-se na extração, processamento em memória e persistência de dados do sistema Tarrafa. Primeiramente, os dados relacionais do Currículo Lattes são lidos via JDBC em paralelo. Na segunda fase, esses dados são transformados em abstrações de grafos (Vértices e Arestas) e distribuídos pelos nós do \textit{cluster Apache Spark}. Na terceira fase, algoritmos de análise de redes são aplicados para identificar padrões de colaboração. Por fim, as métricas geradas são consolidadas e salvas de volta no banco de dados relacional PostgreSQL.

\begin{figure}[H]
    \centering
    % \includegraphics[scale=0.8]{img/sua_imagem_metodologia.png}
    \caption{Ilustração da metodologia proposta para análise distribuída de grafos de colaboração.}
    \label{fig:metodologia}
    \legend{Fonte: Elaborado pelo autor (2026).}
\end{figure}

\section{Materiais}
\label{sec:materiais}

Para o desenvolvimento e avaliação da arquitetura, foram utilizados dados baseados no sistema Tarrafa, que consolida informações extraídas da Plataforma Lattes. O conjunto de dados caracteriza-se por um alto volume de nós (pesquisadores, instituições) e arestas (coautorias em artigos, orientações, projetos conjuntos), formando uma rede de alta complexidade.

Os experimentos serão conduzidos em máquinas com suporte a conteinerização (Docker), respeitando a disponibilidade de recursos. A implementação do pipeline será realizada em \textit{Python} (PySpark), em virtude de sua forte integração com o ecossistema de dados e ferramentas de análise. O banco de dados escolhido foi o PostgreSQL versão 14+.

A avaliação experimental da arquitetura proposta compreende as seguintes análises:

\begin{itemize}
    \item \textbf{Análise de desempenho e I/O:} Avaliação do tempo de leitura e gravação no banco de dados relacional sob diferentes estratégias de particionamento JDBC no Spark.
    
    \item \textbf{Análise de escalabilidade do Grafo:} Comparação do tempo de cálculo de métricas de rede (como graus de separação e conexões conexas) ao variar o número de nós \textit{Workers} e a alocação de memória (\textit{executor memory}).
\end{itemize}

\section{Atividades}
\label{sec:atividades}

A fim de alcançar os objetivos propostos, as seguintes atividades serão realizadas no desenvolvimento deste trabalho:

\begin{enumerate}
    \item Configuração do Apache Spark, PostgreSQL e Python.
    \item Modelagem do esquema relacional e carga inicial de dados amostrais do sistema Tarrafa.
    \item Implementação dos \textit{scripts} PySpark para extração paralela via JDBC.
    \item Construção dos DataFrames de Vértices e Arestas e instanciação do Grafo.
    \item Execução de algoritmos de análise topológica sobre a rede de colaboração.
    \item Implementação do \textit{pipeline} de \textit{upsert} dos resultados para o PostgreSQL.
    \item Realizar experimentos de escalabilidade e coleta de métricas de desempenho.
    \item Análise dos resultados de \textit{shuffle} de rede e consumo de memória.
    \item Documentação dos códigos e resultados obtidos.
    \item Finalização da monografia.
\end{enumerate}

\section{Cronograma}
\label{sec:cronograma}

O cronograma de atividades para o desenvolvimento deste trabalho é apresentado na Tabela \ref{tab:cronograma}, considerando um período de execução de 6 meses.

\begin{table}[H]
\centering
\caption{Cronograma de atividades para o desenvolvimento do trabalho.}
\label{tab:cronograma}
\begin{tabular}{|c|c|c|c|c|c|c|}
\hline
\textbf{Atividade} & \textbf{Set} & \textbf{Out} & \textbf{Nov} & \textbf{Dez} & \textbf{Jan} & \textbf{Fev} \\
\hline
1 & X & X &   &   &   &   \\
\hline
2 & X & X & X &   &   &   \\
\hline
3 &   & X  & X &  &   &   \\
\hline
4 &   &   & X & X &   &   \\
\hline
5 &   &   & X  & X &  &   \\
\hline
6 &   &   &   & X & X &  \\
\hline
7 &   &   &   &  & X & X \\
\hline
8 &   &   &   &   & X & X \\
\hline
9 &   &   &   &   &   & X \\
\hline
10&   &   &   &   &   & X \\
\hline
\end{tabular}
\end{table}